\ifdefined\gkver\else\def\gkver{student}\fi
\documentclass[\gkver]{gaokaozhenti}
\begin{document}

\chapter{2009年广东理科基础}
%% paperType: 理科基础物理部分

\begin{enumerate}
\item
%% number: 1
%% paperName: 2009年广东理科基础第 1 题 2 分
%% typeId: 1
%% type: 单选题
%% chapter: 其他
%% point: 物理学史
%% method: 无
%% score: 2
%% degree: 950
%% duplicateId: 0
%% body:
发现通电导线周围存在磁场的科学家是\xzanswer{D}

\fourchoices[answer=D]
{洛伦兹}
{库仑}
{法拉第}
{奥斯特}

%% answer: D
%% memo:
\memoanswer{
发现电流的磁效应的科学家是丹麦的奥斯特。而法拉第是发现了电磁感应现象。
}

\item
%% number: 2
%% paperName: 2009年广东理科基础第 2 题 2 分
%% typeId: 1
%% type: 单选题
%% chapter: 直线运动
%% point: 质点
%% method: 无
%% score: 2
%% degree: 950
%% duplicateId: 0
%% body:
做下列运动的物体，能当作质点处理的是\xzanswer{D}

\fourchoices[answer=D]
{自转中的地球}
{旋转中的风力发电机叶片}
{在冰面上旋转的花样滑冰运动员}
{匀速直线运动的火车}

%% answer: D
%% memo:
\memoanswer{
A、B、C 选项中对所研究的问题，物体各部分的运动情况不一样所以不能看作质点，D 可以。
}

\item
%% number: 3
%% paperName: 2009年广东理科基础第 3 题 2 分
%% typeId: 1
%% type: 单选题
%% chapter: 直线运动
%% point: 运动图象
%% method: 无
%% score: 2
%% degree: 850
%% duplicateId: 0
%% body:
如图所示是甲、乙两物体做直线运动的$v-t$图象。下列表述正确的是\xzanswer{A}

\onepicture[w=8.0cm]{figs/03.png}

\fourchoices[answer=A]
{乙做匀加速直线运动}
{$0\sim1\Us$ 内甲和乙的位移相等}
{甲和乙的加速度方向相同}
{甲的加速度比乙的小}

%% answer: A
%% memo:
\memoanswer{
甲、乙两物体在速度图象里的图形都是倾斜的直线表明两物体都是匀变速直线，乙是匀加速，甲是匀减速，加速度方向不同 A 对 C 错。根据在速度图象里面积表示位移的方法可知在 $0\sim1\Us$ 内甲通过的位移大于乙通过的位移，B 错。根据斜率表示加速度可知甲的加速度大于乙的加速度，D 错。
}

\item
%% number: 4
%% paperName: 2009年广东理科基础第 4 题 2 分
%% typeId: 1
%% type: 单选题
%% chapter: 运动和力
%% point: 牛顿第二定律
%% method: 无
%% score: 2
%% degree: 850
%% duplicateId: 0
%% body:
建筑工人用如图所示的定滑轮装置运送建筑材料。质量为 $70.0\Ukg$ 的工人站在地面上，通过定滑轮将 $20.0\Ukg$ 的建筑材料以 $0.5\Umsq$ 的加速度拉升，忽略绳子和定滑轮的质量及定滑轮的摩擦，则工人对地面的压力大小为（$g$ 取 $10\Umsq$。）\xzanswer{B}

\onepicture[w=4.1cm]{figs/04.png}

\fourchoices[answer=B]
{$510\UN$}
{$490\UN$}
{$890\UN$}
{$910\UN$}

%% answer: B
%% memo:
\memoanswer{
对建筑材料进行受力分析根据牛顿第二定律有 $F-mg=ma$，得绳子的拉力大小等于 $F=210\UN$，然后再对人受力分析由平衡的知识得 $Mg=F+F_{N}$，得 $F_{N}=490\UN$，根据牛顿第三定律可知人对地面间的压力为 $490\UN$，B 对。
}

\item
%% number: 5
%% paperName: 2009年广东理科基础第 5 题 2 分
%% typeId: 1
%% type: 单选题
%% chapter: 电路
%% point: 电阻定律
%% method: 无
%% score: 2
%% degree: 850
%% duplicateId: 0
%% body:
导体的电阻是导体本身的一种性质，对于同种材料的导体，下列表述正确的是\xzanswer{A}

\fourchoices[answer=A]
{横截面积一定，电阻与导体的长度成正比}
{长度一定，电阻与导体的横截面积成正比}
{电压一定，电阻与通过导体的电流成正比}
{电流一定，电阻与导体两端的电压成反比}

%% answer: A
%% memo:
\memoanswer{
对于同种材料的物体，电阻率是个定值，根据电阻定律 $R=\rho\frac{l}{s}$ 可知 A 正确。
}

\item
%% number: 6
%% paperName: 2009年广东理科基础第 6 题 2 分
%% typeId: 1
%% type: 单选题
%% chapter: 直线运动
%% point: 运动的合成
%% method: 无
%% score: 2
%% degree: 850
%% duplicateId: 0
%% body:
船在静水中的航速为$v_{1}$，水流的速度为$v_{2}$。为使船行驶到河正对岸的码头，则$v_{1}$相对$v_{2}$的方向应为\xzanswer{C}

\fourchoices[answer=C, ispicture=true, wA=2.0cm, wB=1.8cm, wC=3.0cm, wD=3.0cm]
{figs/06a.png}
{figs/06b.png}
{figs/06c.png}
{figs/06d.png}

%% answer: C
%% memo:
\memoanswer{
根据运动的合成与分解的知识，可知要使船垂直达到对岸即要船的合速度指向对岸。根据平行四边行定则，C 能。
}

\item
%% number: 7
%% paperName: 2009年广东理科基础第 7 题 2 分
%% typeId: 1
%% type: 单选题
%% chapter: 曲线运动
%% point: 平抛运动
%% method: 无
%% score: 2
%% degree: 850
%% duplicateId: 0
%% body:
滑雪运动员以$20\Ums$的速度从一平台水平飞出，落地点与飞出点的高度差 $3.2\Um$。不计空气阻力，$g$ 取 $10\Umsq$。运动员飞过的水平距离为$s$，所用时间为$t$，则下列结果正确的是\xzanswer{B}

\fourchoices[answer=B]
{$s=16\Um$，$t=0.50\Us$}
{$s=16\Um$，$t=0.80\Us$}
{$s=20\Um$，$t=0.50\Us$}
{$s=20\Um$，$t=0.80\Us$}

%% answer: B
%% memo:
\memoanswer{
做平抛运动的物体运动时间由高度决定，根据竖直方向做自由落体运动得 $t=\sqrt{\frac{2h}{g}}=0.80\Us$，根据水平方向做匀速直线运动可知 $s=v_{0}t=20\times0.80=16\Um$，B 正确。
}

\item
%% number: 8
%% paperName: 2009年广东理科基础第 8 题 2 分
%% typeId: 1
%% type: 单选题
%% chapter: 功和能
%% point: 功
%% method: 无
%% score: 2
%% degree: 850
%% duplicateId: 0
%% body:
游乐场中的一种滑梯如图所示。小朋友从轨道顶端由静止开始下滑，沿水平轨道滑动了一段距离后停下来，则\xzanswer{D}

\onepicture[w=10.1cm]{figs/08.png}

\fourchoices[answer=D]
{下滑过程中支持力对小朋友做功}
{下滑过程中小朋友的重力势能增加}
{整个运动过程中小朋友的机械能守恒}
{在水平面滑动过程中摩擦力对小朋友做负功}

%% answer: D
%% memo:
\memoanswer{
在滑动的过程中，人受三个力重力做正功，势能降低 B 错，支持力不做功，摩擦力做负功，所以机械能不守恒，A、C 皆错。D 正确。
}

\item
%% number: 9
%% paperName: 2009年广东理科基础第 9 题 2 分
%% typeId: 1
%% type: 单选题
%% chapter: 功和能
%% point: 功能中的图象
%% method: 无
%% score: 2
%% degree: 850
%% duplicateId: 0
%% body:
物体在合外力作用下做直线运动的$v-t$图象如图所示。下列表述正确的是\xzanswer{A}

\onepicture[w=8.15cm]{figs/09.png}

\fourchoices[answer=A]
{在 $0\sim1\Us$ 内，合外力做正功}
{在 $0\sim2\Us$ 内，合外力总是做负功}
{在 $1\sim2\Us$ 内，合外力不做功}
{在 $0\sim3\Us$ 内，合外力总是做正功}

%% answer: A
%% memo:
\memoanswer{
根据物体的速度图象可知，物体 $0\sim1\Us$ 内做匀加速合外力做正功，A 正确。$1\sim3\Us$ 内做匀减速合外力做负功。根据动能定理 $0\sim3\Us$ 内，$1\sim2\Us$ 内合外力做功为零。
}

\item
%% number: 10
%% paperName: 2009年广东理科基础第 10 题 2 分
%% typeId: 1
%% type: 单选题
%% chapter: 曲线运动
%% point: 人造地球卫星
%% method: 无
%% score: 2
%% degree: 950
%% duplicateId: 0
%% body:
关于地球的第一宇宙速度，下列表述正确的是\xzanswer{A}

\fourchoices[answer=A]
{第一宇宙速度又叫环绕速度}
{第一宇宙速度又叫脱离速度}
{第一宇宙速度跟地球的质量无关}
{第一宇宙速度跟地球的半径无关}

%% answer: A
%% memo:
\memoanswer{
第一宇宙速度又叫环绕速度 A 对，B 错。根据定义有 $G\frac{mM}{R^{2}}=m\frac{v^{2}}{R}$ 可知与地球的质量和半径有关。C、D 错。
}

\item
%% number: 11
%% paperName: 2009年广东理科基础第 11 题 2 分
%% typeId: 1
%% type: 单选题
%% chapter: 曲线运动
%% point: 人造地球卫星
%% method: 无
%% score: 2
%% degree: 850
%% duplicateId: 0
%% body:
宇宙飞船在半径为$R_{1}$的轨道上运行，变轨后的半径为$R_{2}$，$R_{1}>R_{2}$。宇宙飞船绕地球做匀速圆周运动，则变轨后宇宙飞船的\xzanswer{D}

\fourchoices[answer=D]
{线速度变小}
{角速度变小}
{周期变大}
{向心加速度变大}

%% answer: D
%% memo:
\memoanswer{
根据 $G\frac{mM}{r^{2}}=m\frac{v^{2}}{r}=m\omega^{2}r=m\frac{4\pi^{2}r}{T^{2}}=ma_{\text{向}}$ 得 $v=\sqrt{\frac{GM}{r}}$，可知变轨后飞船的线速度变大，A 错。角速度变大 B 错。周期变小 C 错。向心加速度在增大 D 正确。
}

\item
%% number: 12
%% paperName: 2009年广东理科基础第 12 题 2 分
%% typeId: 1
%% type: 单选题
%% chapter: 电场
%% point: 电场线
%% method: 无
%% score: 2
%% degree: 850
%% duplicateId: 0
%% body:
关于同一电场的电场线，下列表述正确的是\xzanswer{C}

\fourchoices[answer=C]
{电场线是客观存在的}
{电场线越密，电场强度越小}
{沿着电场线方向，电势越来越低}
{电荷在沿电场线方向移动时，电势能减小}

%% answer: C
%% memo:
\memoanswer{
电场是客观存在的，而电场线是假想的，A 错。电场线越密的地方电场越大 B 错。沿着电场线的方向电势逐渐降低 C 对，负电荷沿着电场线方向移动时电场力做负功电势能增加 D 错。
}

\item
%% number: 13
%% paperName: 2009年广东理科基础第 13 题 2 分
%% typeId: 1
%% type: 单选题
%% chapter: 磁场
%% point: 洛伦兹力
%% method: 无
%% score: 2
%% degree: 850
%% duplicateId: 0
%% body:
带电粒子垂直匀强磁场方向运动时，会受到洛伦兹力的作用。下列表述正确的是\xzanswer{B}

\fourchoices[answer=B]
{洛伦兹力对带电粒子做功}
{洛伦兹力不改变带电粒子的动能}
{洛伦兹力的大小与速度无关}
{洛伦兹力不改变带电粒子的速度方向}

%% answer: B
%% memo:
\memoanswer{
根据洛伦兹力的特点，洛伦兹力对带电粒子不做功，A 错，B 对。根据 $F=qvB$，可知大小与速度有关，洛伦兹力的效果就是改变物体的运动方向，不改变速度的大小。
}

\item
%% number: 14
%% paperName: 2009年广东理科基础第 14 题 2 分
%% typeId: 1
%% type: 单选题
%% chapter: 电路
%% point: 动态分析
%% method: 无
%% score: 2
%% degree: 850
%% duplicateId: 0
%% body:
如图所示是一实验电路图。在滑动触头由 a 端滑向 b 端的过程中，下列表述正确的是\xzanswer{A}

\onepicture[w=7.8cm]{figs/14.png}

\fourchoices[answer=A]
{路端电压变小}
{电流表的示数变大}
{电源内阻消耗的功率变小}
{电路的总电阻变大}

%% answer: A
%% memo:
\memoanswer{
当滑片向 b 端滑动时，接入电路中的电阻减少，使得总电阻减小 D 错。根据 $I=\frac{E}{R_{\text{总}}}$，可知总电流在增加，根据闭合电路中的欧姆定律有 $E=Ir+U_{\text{外}}$，可知路端电压在减小，A 对。流过电流表的示数为 $I=\frac{U_{\text{外}}}{R_{3}}$，可知电流在减小，B 错。根据 $P=I^{2}r$，可知内阻消耗的功率在增大，C 错。
}

\item
%% number: 15
%% paperName: 2009年广东理科基础第 15 题 2 分
%% typeId: 1
%% type: 单选题
%% chapter: 运动和力
%% point: 牛顿第二定律
%% method: 无
%% score: 2
%% degree: 750
%% duplicateId: 0
%% body:
搬运工人沿粗糙斜面把一个物体拉上卡车，当力沿斜面向上，大小为$F$时，物体的加速度为$a_{1}$；若保持力的方向不变，大小变为 $2F$ 时，物体的加速度为$a_{2}$，则\xzanswer{D}

\fourchoices[answer=D]
{$a_{1}=a_{2}$}
{$a_{1}<a_{2}<2a_{1}$}
{$a_{2}=2a_{1}$}
{$a_{2}>2a_{1}$}

%% answer: D
%% memo:
\memoanswer{
当为 $F$ 时有 $a_{1}=\frac{F-f}{m}$，当为 $2F$ 时有 $a_{2}=\frac{2F-f}{m}=\frac{2F-2f+f}{m}=2a_{1}+\frac{f}{m}$，可知 $a_{2}>2a_{1}$，D 对。
}

\item
%% number: 16
%% paperName: 2009年广东理科基础第 16 题 2 分
%% typeId: 1
%% type: 单选题
%% chapter: 电场
%% point: 带电粒子的轨迹类问题
%% method: 无
%% score: 2
%% degree: 650
%% duplicateId: 0
%% body:
如图所示，一带负电粒子以某速度进入水平向右的匀强电场中，在电场力作用下形成图中所示的运动轨迹。M 和 N 是轨迹上的两点，其中 M 点在轨迹的最右点。不计重力，下列表述正确的是\xzanswer{C}

\onepicture[w=5.3cm]{figs/16.png}

\fourchoices[answer=C]
{粒子在 M 点的速率最大}
{粒子所受电场力沿电场方向}
{粒子在电场中的加速度不变}
{粒子在电场中的电势能始终在增加}

%% answer: C
%% memo:
\memoanswer{
根据做曲线运动物体的受力特点合力指向轨迹的凹一侧，再结合电场力的特点可知粒子带负电，即受到的电场力方向与电场线方向相反，B 错。从 N 到 M 电场力做负功，减速，电势能在增加，当达到 M 点后电场力做正功加速电势能在减小则在 M 点的速度最小 A 错，D 错。在整个过程中只受电场力根据牛顿第二定律加速度不变。
}

\item
%% number: 17
%% paperName: 2009年广东理科基础第 17 题 2 分
%% typeId: 1
%% type: 单选题
%% chapter: 力物体的平衡
%% point: 弹力
%% method: 无
%% score: 2
%% degree: 650
%% duplicateId: 0
%% body:
一个实验小组在“探究弹力和弹簧伸长的关系”的实验中，使用两条不同的轻质弹簧 a 和 b，得到弹力与弹簧长度的图象如图所示。下列表述正确的是\xzanswer{B}

\onepicture[w=7.1cm]{figs/17.png}

\fourchoices[answer=B]
{a 的原长比 b 的长}
{a 的劲度系数比 b 的大}
{a 的劲度系数比 b 的小}
{测得的弹力与弹簧的长度成正比}

%% answer: B
%% memo:
\memoanswer{
图象中的斜率表示劲度系数，可知 a 的劲度系数比 b 的大，B 正确。与 $l$ 的截距表示原长则 a 的原长比 b 的短，A 错。
}

\item
%% number: 18
%% paperName: 2009年广东理科基础第 18 题 2 分
%% typeId: 1
%% type: 单选题
%% chapter: 直线运动
%% point: 打点计时器公式
%% method: 无
%% score: 2
%% degree: 750
%% duplicateId: 0
%% body:
“研究匀变速直线运动”的实验中，使用电磁式打点计时器（所用交流电的频率为$50\UHz$），得到如图所示的纸带。图中的点为计数点，相邻两计数点间还有四个点未画出来，下列表述正确的是\xzanswer{C}

\onepicture[w=16.5cm]{figs/18.png}

\fourchoices[answer=C]
{实验时应先放开纸带再接通电源}
{（$s_{6}-s_{1}$）等于（$s_{2}-s_{1}$）的 6 倍}
{从纸带可求出计数点 B 对应的速率}
{相邻两个计数点间的时间间隔为$0.02\Us$}

%% answer: C
%% memo:
\memoanswer{
在“研究匀变速直线运动”的实验中，实验时应先接通电源再放开纸带，A 错。根据相等的时间间隔内通过的位移有 $x_{M}-x_{N}=(M-N)at^{2}$，可知 $(s_{6}-s_{1})$ 等于 $(s_{2}-s_{1})$ 的 5 倍，B 错。根据 B 点为 A 与 C 的中间时刻点有 $v_{B}=\frac{x_{AC}}{2t}$，C 对。由于相邻的计数点之间还有 4 个点没有画出，所以时间间隔为 $0.1\Us$，D 错。
}

\end{enumerate}

\end{document}
